% English translation of checked_pages/esq_119_003.tex; source: cote 119, site 3, PDF page 4, typescript page 1.
% AI-translated by GPT-6 Astra (OpenAI, ChatGPT Pro mode).
\scanpage{119/3}
\mspage{1}

SKETCH OF A PROGRAM

by A. Grothendieck

1. As the present circumstances make the prospects of teaching at research level
at the University increasingly illusory for me, I have resolved to apply for
admission to the CNRS,\tnote{TN-003-01}{CNRS: \href{https://aissai.cnrs.fr/en/legal-notice/}{Centre national de la recherche scientifique}, the French National Center for Scientific Research.}
so that I may devote my energy to developing work and perspectives
that, it is becoming clear, no student (nor even, it seems,
any fellow mathematician) will be found to develop in my
stead.

In place of a ``Qualifications and Works'' document,\tnote{TN-003-02}{``Titres et Travaux'' is the French designation of the document listing qualifications and scholarly work.}
following this text will be found the full reproduction of a sketch, arranged by themes,
of what I regarded as my main mathematical contributions at the time
of writing that report, in 1972. It also contains
a list of articles published by that date. I have published no scientific
articles since 1970. In the lines that follow,
I propose to give at least a glimpse of some principal themes
of my mathematical reflections since then. Over the years these reflections
have taken material form in two bulky boxes of
manuscript notes, probably difficult for anyone other than
myself to decipher, which, after successive stages of decantation,
are perhaps awaiting their moment for an overall write-up, at least
a provisional one, intended for the mathematical community. The
term ``write-up'' here is somewhat inappropriate,\tnote{TN-003-03}{``Write-up'' renders ``rédaction''. The source itself immediately explains the continuing work of discovery meant by this word.}
since it is much more a matter of developing the many ideas and visions begun
over these last twelve years, making them more precise and exploring them
more deeply, with all the unforeseen turns that constantly
accompany this kind of work - a work of discovery, then, and
not of compiling notes piously accumulated. And I fully intend,
in writing the ``Mathematical Reflections'' begun in February
1983,\tnote{TN-003-04}{The French title is ``Réflexions Mathématiques''.}
to let the course of thought that probes and discovers appear
clearly through the pages, often groping in the half-light,
with sudden openings of light when some stubborn

\margin{3}

